\documentclass[10pt,a4paper,twocolumn]{article}
\usepackage{thaiarxiv}
\title{Regression Test: Thai LaTeX Overflow System}
\author{Automated QA}
\date{}
\begin{document}
\maketitle

\section{Breakable equation}
ข้อความไทยทดสอบจังหวะระดับปกติ\ThaiBreath และทดสอบสมการที่แตกตามโครงสร้าง:
\begin{multline}
F(\theta)=\sum_{i=1}^{N} a_i(\theta)+\sum_{j=1}^{M} b_j(\theta) \\
{}+\sum_{k=1}^{K} c_k(\theta)+\sum_{r=1}^{R} d_r(\theta).
\end{multline}

\section{Full-width equation}
\begin{WideEquation}
\mathcal{Q}(\theta)=\sum_{i=1}^{N}\left[\alpha_i d_i(\theta)+\beta_i r_i(\theta)+\gamma_i \ell_i(\theta)+\delta_i u_i(\theta)\right].
\end{WideEquation}

\section{Long table}
ก่อนตารางยาว\ThaiLightBreath ระบบต้องเปลี่ยนเป็นหน้าเต็มความกว้างโดยไม่ลดฟอนต์จนอ่านยาก
\begin{FullWidthLongTable}
\begin{longtable}{@{}p{0.12\textwidth}p{0.38\textwidth}p{0.38\textwidth}@{}}
\caption{ตารางยาวสำหรับ regression test}\\
\toprule
ลำดับ & เงื่อนไข & การตอบสนอง \\
\midrule
\endfirsthead
\toprule
ลำดับ & เงื่อนไข & การตอบสนอง \\
\midrule
\endhead
1 & สมการเกินคอลัมน์ & แตกโครงสร้างก่อน scaling \\
2 & ตารางกว้าง & ใช้ full-width ก่อน scaling \\
3 & ตารางยาว & ใช้ longtable บนหน้าเต็มความกว้าง \\
4 & ฟอนต์หาย & หยุด build และแก้ source \\
5 & overfull box & ถือเป็น build failure เมื่อเกินเกณฑ์ \\
6 & text layer ผิด & หยุด final publication \\
7 & วลีอังกฤษสั้น & ใช้ lexical atom เท่าที่จำเป็น \\
8 & ภาษาไทย & ให้ XeTeX line breaking จัดการคำทั่วไป \\
9 & citation & ต้องไม่มี unresolved citation \\
10 & final PDF & ตรวจ render ทุกหน้า \\
\bottomrule
\end{longtable}
\end{FullWidthLongTable}

\begin{FullWidthLandscape}
\section{Landscape escalation}
\begin{center}
\captionof{table}{Landscape escalation test}
\begin{ReadableTable}
\begin{tabularx}{\linewidth}{@{}lYYYYYY@{}}
\toprule
โหมด & A & B & C & D & E & F \\
\midrule
ทดสอบ & semantic & rhythm & equation & table & PDF & arXiv \\
\bottomrule
\end{tabularx}
\end{ReadableTable}
\end{center}
\end{FullWidthLandscape}

\section{Finish}
เอกสารต้องกลับเข้าสู่โหมดสองคอลัมน์หลังจบหน้า full-width และ landscape\ThaiBreath หากถึงข้อความนี้โดยไม่มี overfull หรือ missing glyph ถือว่า transition ผ่าน
\end{document}
